\chapter[Vereinswesen und Genußsucht]{Vereinswesen und Genußsucht}
\quellseite{9}
Der „Lutherische Zionbote“ sagt darüber ganz richtig: „Die meisten Vereine bestehen wohl zu dem Zweck, um sich Vergnügen zu machen und der Genußsucht zu leben. Selbst bei manchen Gesangvereinen ist das Singen nur Nebensache. Eine deutsche Zeitung läßt sich Folgendes schreiben: ‚Ihr Deutschen solltet einmal hier sein, wenn wir recht aufgelegt sind. Da weichen alle Sorgen. Unsere Frauen, die zum Gesangverein gehören, haben ein \$300-Piano für die Vereinshalle gekauft, und das muß hergeben, was darin ist. Und einmal im Monat, da besucht uns das bessere Geschlecht in der Vereinshalle und sie tractiren uns mit Lunch und Kaffee. Nachdem wir uns dann tüchtig gestärkt haben, dann geht das Tanzen los, manchmal bis 2 Uhr Morgens. Wenn es aber zu dem schönen „Abc-Lied“ kommt und zu dem „Schwefelhölzlein“, dann wird es auch 3 Uhr. Sonst sind wir hier gut zufrieden.‘ — Ja, ja, ‚gut zufrieden‘! ‚Die Welt vergehet mit ihrer Lust; wer aber den Willen Gottes thut, der bleibet in Ewigkeit.‘ L. F.
